\chapter[21]{Perelman's Pseudolocality Theorem}
\label{ch:chow-iii-perelmans-pseudolocality}

Throughout this chapter, \((\mathcal M^n,g(t))\) denotes a complete Ricci flow
when completeness is stated explicitly.  We write \(\Rm\) and \(\Rc\) for the
Riemann and Ricci tensors, \(\Area\) and \(\Vol\) for hypersurface area and
volume, and \(c_n=n^n\omega_n\) for the Euclidean isoperimetric constant.

\section{Statement and interpretation of pseudolocality}

\begin{proposition}[Euclidean isoperimetric inequality]
\label{prop:chow-iii-euclidean-isoperimetric}
\source{chow-et-al-ricci-flow-part-iii:page-0180}
\dcref{ch21:21.1}
\group{Pseudolocality hypotheses}

For every regular domain \(\Omega\subset\mathbb E^n\),
\[
  \Area(\partial\Omega)^n\ge n^n\omega_n\,\Vol(\Omega)^{n-1}=c_n\Vol(\Omega)^{n-1}.
\]
\end{proposition}

\begin{theorem}[Perelman's pseudolocality]
\label{thm:chow-iii-perelman-pseudolocality}
\source{chow-et-al-ricci-flow-part-iii:page-0180}
\uses{prop:chow-iii-euclidean-isoperimetric}
\dcref{ch21:21.2}
\group{Pseudolocality theorem}


For every \(\alpha>0\) and \(n\ge2\), there are \(\delta,\epsilon_0>0\),
depending only on \(\alpha,n\), such that the following holds.  If
\(t\in[0,(\epsilon r_0)^2]\), \(\epsilon\le\epsilon_0\), and \(g(t)\) is a
complete bounded-curvature Ricci flow with
\[
\Rm(x,0)\ge-r_0^{-2}\quad\text{on }B_{g(0)}(x_0,r_0)
\]
and every regular \(\Omega\subset B_{g(0)}(x_0,r_0)\) satisfies
\[
\Area_{g(0)}(\partial\Omega)^n\ge(1-\delta)c_n\Vol_{g(0)}(\Omega)^{n-1},
\]
then
\[
|\Rm|(x,t)\le\frac{\alpha}{t}+\frac1{(\epsilon_0r_0)^2}
\]
whenever \(d_{g(t)}(x,x_0)<\epsilon_0r_0\) and
\(0<t\le(\epsilon r_0)^2\).
\end{theorem}

\begin{remark}[Size of the time interval]
\label{rem:chow-iii-pseudolocality-scales}
\source{chow-et-al-ricci-flow-part-iii:page-0180}
\dcref{ch21:21.3}
\group{Pseudolocality theorem}

The time interval uses \(\epsilon\), while the spatial radius and additive
curvature term use \(\epsilon_0\).
\end{remark}

\begin{corollary}[Interior curvature estimate]
\label{cor:chow-iii-interior-curvature-estimate}
\source{chow-et-al-ricci-flow-part-iii:page-0181}
\uses{thm:chow-iii-perelman-pseudolocality}
\dcref{ch21:21.4}
\group{Pseudolocality theorem}


If the hypotheses of Theorem~\ref{thm:chow-iii-perelman-pseudolocality} hold,
then, for every \(\eta\in(0,1)\),
\[
|\Rm|\le\frac\alpha t+\frac1{(\epsilon_0\eta r_0)^2}
\]
on \(B_{g(0)}(x_0,(1-\eta)r_0)\times(0,(\epsilon\eta r_0)^2]\).
\end{corollary}

\begin{remark}[Scale-invariance of the statement]
\label{rem:chow-iii-pseudolocality-scale-invariance}
\source{chow-et-al-ricci-flow-part-iii:page-0181}
\dcref{ch21:21.5}
\group{Pseudolocality theorem}

Under \(\widetilde g(t)=Cg(C^{-1}t)\) and \(\widetilde r_0=C^{1/2}r_0\),
the curvature lower bound, almost-Euclidean isoperimetric condition, and
the conclusion of Theorem~\ref{thm:chow-iii-perelman-pseudolocality} are
unchanged.  In particular, \(|\Rm_{\widetilde g(t)}|=C^{-1}|\Rm_{g(C^{-1}t)}|\).
\end{remark}

\begin{example}[Topping's incomplete solution]
\label{ex:chow-iii-topping-incomplete-solution}
\source{chow-et-al-ricci-flow-part-iii:page-0184,chow-et-al-ricci-flow-part-iii:page-0185}
\dcref{ch21:21.6}
\group{Completeness is essential}

Let \(g_0^r\) be a nonnegatively curved metric on a capped thin cylinder
\(S^1(r)\times[-1,1]\), with
\(4\pi r\le\Area_{g_0^r}(\Sigma)\le(22/5)\pi r\).  Its Ricci flow on
\(S^2\) becomes extinct at \(T^r=\Area_{g_0^r}(\Sigma)/(8\pi)\le11r/20\).
Pulling the flow back through the local covering
\(\phi:(-3/5,3/5)^2\to S^1(r)\times[-1,1]\) gives an incomplete flat-at-time-zero
solution \(g_{\mathcal M}^r(t)\) whose scalar curvature tends uniformly to
infinity as \(t\nearrow T^r\), although every compact initial ball is Euclidean.
\end{example}

\begin{theorem}[Topping]
\label{thm:chow-iii-topping-counterexample}
\source{chow-et-al-ricci-flow-part-iii:page-0185}
\uses{ex:chow-iii-topping-incomplete-solution}
\dcref{ch21:21.7}
\group{Completeness is essential}


There are counterexamples to Theorem~\ref{thm:chow-iii-perelman-pseudolocality}
if completeness of the evolving solution is omitted.  Indeed, choose
\(x_0=(0,0)\), \(r_0=1/2\), and let \(r\to0\) so that \(T^r<
(\epsilon_0r_0)^2\); the claimed estimate would uniformly bound the curvature
up to extinction, contradicting \(\inf R_{g_{\mathcal M}^r(t)}\to\infty\).
\end{theorem}

\begin{remark}[Higher-dimensional incomplete counterexamples]
\label{rem:chow-iii-topping-products}
\source{chow-et-al-ricci-flow-part-iii:page-0186}
\dcref{ch21:21.8}
\group{Completeness is essential}

Taking products of Topping's surface flow with flat tori gives counterexamples
in every dimension at least two.
\end{remark}

\section{Proof by contradiction and point picking}

\begin{theorem}[\(r_0=1\) case of pseudolocality]
\label{thm:chow-iii-pseudolocality-unit-scale}
\source{chow-et-al-ricci-flow-part-iii:page-0188}
\dcref{ch21:21.9}
\group{Contradiction setup}

For every \(\alpha>0\) and \(n\ge2\), there are \(\delta,\epsilon_0>0\) such
that a complete bounded-curvature flow on \([0,\epsilon^2]\), \(\epsilon\le
\epsilon_0\), satisfying \(R(\cdot,0)\ge-1\) and the \((1-\delta)\)-Euclidean
isoperimetric inequality in \(B_{g(0)}(x_0,1)\), obeys
\[
|\Rm|(x,t)\le\frac\alpha t+\frac1{\epsilon_0^2}
\]
for \(d_{g(t)}(x,x_0)<\epsilon_0\) and \(0<t\le\epsilon^2\).
\end{theorem}

\begin{remark}[Monotonicity of \(\delta\), \(\epsilon_0\), and \(\alpha\)]
\label{rem:chow-iii-pseudolocality-parameter-monotonicity}
\source{chow-et-al-ricci-flow-part-iii:page-0188}
\uses{thm:chow-iii-pseudolocality-unit-scale}
\dcref{ch21:21.10}
\group{Contradiction setup}


If \((\delta,\epsilon_0)\) works for \(\alpha\), then any smaller positive
\(\delta',\epsilon_0'\) works for every \(\alpha'\ge\alpha\).
\end{remark}

\begin{definition}[\(\alpha\)-large curvature points]
\label{def:chow-iii-alpha-large-curvature}
\source{chow-et-al-ricci-flow-part-iii:page-0189}
\dcref{ch21:21.11}
\group{Contradiction setup}

For a flow on \([0,\epsilon^2]\), set
\[
\mathcal M_\alpha=\{(x,t):0<t\le\epsilon^2,\ |\Rm|(x,t)>\alpha/t\}.
\]
\end{definition}

\begin{theorem}[Counterstatement A (bad sequence)]
\label{ctr:chow-iii-pseudolocality-counterstatement}
\source{chow-et-al-ricci-flow-part-iii:page-0189}
\uses{def:chow-iii-alpha-large-curvature}
\dcref{ch21:21.12--21.15}
\group{Contradiction setup}


If Theorem~\ref{thm:chow-iii-pseudolocality-unit-scale} fails, there are
\(\delta_i,\epsilon_{0i},\epsilon_i\to0\), complete pointed flows
\((\mathcal M_i,g_i(t),x_{0i})\), and \((x_i,t_i)\in\mathcal M_{i\alpha}\) with
\(d_{g_i(t_i)}(x_i,x_{0i})<\epsilon_i\) and
\[
|\Rm_{g_i}|(x_i,t_i)>\alpha/t_i+\epsilon_{0i}^{-2},
\]
while \(R_{g_i}(\cdot,0)\ge-1\) and the initial isoperimetric constants are
\((1-\delta_i)c_n\).
\end{theorem}

\begin{lemma}[An adjustment for the sequences \(\epsilon_i\) and \((x_i,t_i)\)]
\label{lem:chow-iii-adjusted-bad-sequence}
\source{chow-et-al-ricci-flow-part-iii:page-0189,chow-et-al-ricci-flow-part-iii:page-0190}
\uses{ctr:chow-iii-pseudolocality-counterstatement}
\dcref{ch21:21.12}
\group{Point picking}


After shrinking \(\epsilon_i\) and replacing the bad point if necessary, one
may arrange
\[
|\Rm_{g_i}|(x,t)\le\alpha/t+2\epsilon_{0i}^{-2}
\]
whenever \(0<t\le\epsilon_i^2\) and
\(d_{g_i(t)}(x,x_{0i})\le\epsilon_{0i}\), while retaining a bad point.
\end{lemma}

\begin{remark}[Local bound after adjustment]
\label{rem:chow-iii-adjusted-local-bound}
\source{chow-et-al-ricci-flow-part-iii:page-0190}
\uses{lem:chow-iii-adjusted-bad-sequence}
\dcref{ch21:21.13}
\group{Point picking}


The adjustment implies \(|\Rm_{g_i}|(x,t)<2|\Rm_{g_i}|(x_i,t_i)|\) for
\(t\in[t_i/2,\epsilon_i^2]\) and
\(d_{g_i(t)}(x,x_{0i})\le\epsilon_{0i}\).
\end{remark}

\begin{lemma}[Well-chosen large-curvature points]
\label{lem:chow-iii-well-chosen-points}
\source{chow-et-al-ricci-flow-part-iii:page-0191}
\uses{lem:chow-iii-adjusted-bad-sequence}
\dcref{ch21:21.20--21.27}
\group{Point picking}


With \(A_i=(100n\epsilon_{0i})^{-1}\to\infty\), point picking produces
\((\bar x_i,\bar t_i)\in\mathcal M_{i\alpha}\),
\(\bar Q_i=|\Rm|(\bar x_i,\bar t_i)\), with
\(\bar x_i\in B_{g_i(\bar t_i)}(x_{0i},(2A_i+1)\epsilon_{0i})\), and the
backward curvature bound \(|\Rm|\le4\bar Q_i\) on the parabolic neighborhood
specified by (21.25)--(21.26).  Moreover \(\bar t_i>\alpha\bar Q_i^{-1}\).
\end{lemma}

\section{Local entropies near bad points}

\begin{definition}[Adjoint heat kernels and differential Harnack quantities]
\label{def:chow-iii-adjoint-entropy-quantities}
\source{chow-et-al-ricci-flow-part-iii:page-0192}
\uses{lem:chow-iii-well-chosen-points}
\dcref{ch21:21.28--21.32}
\group{Local entropy}


Let \(H_i\) be the minimal positive fundamental solution of
\(\Box_i^*H_i=(-\partial_t-\Delta_{g_i}+R_{g_i})H_i=0\), based at
\((\bar x_i,\bar t_i)\), normalized by \(\int H_i\,d\mu_{g_i(t)}=1\).  Write
\(H_i=(4\pi(\bar t_i-t))^{-n/2}e^{-f_i}\) and define
\[
v_i=\bigl[(\bar t_i-t)(R_{g_i}+2\Delta f_i-|\nabla f_i|^2)+f_i-n\bigr]H_i.
\]
Perelman's differential Harnack estimate gives \(v_i\le0\).
\end{definition}

\begin{theorem}[Claim 3 (uniform negative local entropy)]
\label{clm:chow-iii-uniform-negative-local-entropy}
\source{chow-et-al-ricci-flow-part-iii:page-0193}
\uses{def:chow-iii-adjoint-entropy-quantities}
\dcref{ch21:21.33--21.35}
\group{Local entropy}


After passing to a subsequence, there is \(\beta_1>0\) and
\(\widetilde t_i\in[\bar t_i-\frac\alpha2\bar Q_i^{-1},\bar t_i)\) such that
\[
\int_{B_{g_i(\widetilde t_i)}(\bar x_i,\sqrt{\bar t_i-\widetilde t_i})}
v_i(\widetilde t_i)\,d\mu_{g_i(\widetilde t_i)}\le-\beta_1.
\]
This assertion is invariant under the parabolic dilations (21.34).
\end{theorem}

\begin{lemma}[Noncollapsed limit and entropy rigidity]
\label{lem:chow-iii-noncollapsed-entropy-limit}
\source{chow-et-al-ricci-flow-part-iii:page-0194,chow-et-al-ricci-flow-part-iii:page-0195,chow-et-al-ricci-flow-part-iii:page-0196,chow-et-al-ricci-flow-part-iii:page-0197}
\uses{clm:chow-iii-uniform-negative-local-entropy}
\dcref{ch21:21.36--21.54}
\group{Local entropy}


If \(\bar Q_i^{1/2}\operatorname{inj}_{g_i(\bar t_i)}(\bar x_i)\) is bounded below, the
rescaled flows converge to a complete bounded-curvature limit
\((\bar{\mathcal M}_\infty,\bar g_\infty(t),\bar x_\infty)\) on
\((-\alpha/2,0]\), with \(|\Rm|(\bar x_\infty,0)=1\).  The rescaled entropy
quantities converge to \(\bar v_\infty\le0\) satisfying
\[
\Box_\infty^*\bar v_\infty
 =2t\left|\Rc_{\bar g_\infty}+\bar\nabla^2\bar f_\infty+\frac1{2t}\bar g_\infty\right|^2\bar H_\infty\le0.
\]
If the local entropy conclusion failed, the strong maximum principle would
force \(\bar v_\infty\equiv0\), making the limit a shrinking gradient soliton
smooth through \(t=0\); bounded curvature then forces it to be flat, contradicting
\(|\Rm|(\bar x_\infty,0)=1\).
\end{lemma}

\begin{lemma}[Collapsed limit and injectivity-radius rigidity]
\label{lem:chow-iii-collapsed-entropy-limit}
\source{chow-et-al-ricci-flow-part-iii:page-0197,chow-et-al-ricci-flow-part-iii:page-0198,chow-et-al-ricci-flow-part-iii:page-0199}
\uses{clm:chow-iii-uniform-negative-local-entropy}
\dcref{ch21:21.55--21.65}
\group{Local entropy}


If \(\bar Q_i^{1/2}\operatorname{inj}_{g_i(\bar t_i)}(\bar x_i)\to0\), rescale instead so that
the injectivity radius at the base point is one.  Hamilton compactness gives a
complete ancient flat limit \((\widehat{\mathcal M}_\infty,\widehat g_\infty(t))\)
with \(\operatorname{inj}_{\widehat g_\infty(0)}(\bar x_\infty)=1\).  The limiting entropy
\(\widehat v_\infty\le0\) satisfies the analogue of the differential Harnack
identity.  Vanishing of its integral on any centered ball would force a flat
shrinking soliton, hence stationary Euclidean space, contradicting injectivity
radius one.  Therefore the negative local entropy bound in Claim~\ref{clm:chow-iii-uniform-negative-local-entropy} also holds in the collapsed case.
\end{lemma}

\section{Contradiction with the logarithmic Sobolev inequality}

\begin{lemma}[Curvature gap for Type III solutions]
\label{lem:chow-iii-type-iii-curvature-gap}
\source{chow-et-al-ricci-flow-part-iii:page-0201}
\dcref{ch21:21.19}
\group{Notes and commentary}

If an eternal solution on a compact manifold satisfies
\(\max|\Rm|(t)\le\alpha/t\) for all sufficiently large \(t\), with
\(\alpha<1/(2(n-1))\), then
\[
\operatorname{diam}(g(t))^2\max|\Rm|(t)\longrightarrow0,
\]
so the metrics become almost flat in the sense of Gromov.
\end{lemma}

\begin{remark}[Flat limits and finite covers]
\label{rem:chow-iii-flat-limit-cover}
\source{chow-et-al-ricci-flow-part-iii:page-0202}
\uses{lem:chow-iii-collapsed-entropy-limit}
\dcref{ch21:21.20}
\group{Notes and commentary}


A complete noncompact flat limit has a finite cover isometric to
\(\mathbb R^k\times T^{n-k}\), and is a flat vector bundle over a closed flat
manifold.  The lifted adjoint heat kernel is the finite deck-transformation sum
of the kernel on the cover, although the nonlinear Harnack quantity does not
decompose as a corresponding sum.
\end{remark}

\begin{theorem}[Completion of the unit-scale proof]
\label{thm:chow-iii-pseudolocality-log-sobolev-contradiction}
\source{chow-et-al-ricci-flow-part-iii:page-0199,chow-et-al-ricci-flow-part-iii:page-0200}
\uses{clm:chow-iii-uniform-negative-local-entropy, lem:chow-iii-noncollapsed-entropy-limit, lem:chow-iii-collapsed-entropy-limit, thm:chow-iii-pseudolocality-unit-scale}
\dcref{ch21:21.66--21.68}
\group{Logarithmic Sobolev contradiction}


The negative local entropy furnished by Claim~\ref{clm:chow-iii-uniform-negative-local-entropy}
propagates to time zero through the localized adjoint heat-kernel estimate.  After
the normalization \(\dot g_i(0)=(2\widetilde t_i)^{-1}g_i(0)\) and
\(\psi_i^2=(2\widetilde t_i)^{n/2}\widetilde H_i\), it yields
\[
\int\bigl(2|\nabla\psi_i|^2-\psi_i^2\log\psi_i^2-s_n\psi_i^2\bigr)
\,d\mu_{\dot g_i(0)}\le-\frac{\beta_1}{2}\int\psi_i^2\,d\mu_{\dot g_i(0)}.
\]
The almost Euclidean isoperimetric inequality and the logarithmic Sobolev
inequality give the opposite lower bound with coefficient
\(s_n+\log(1-\delta_i)\).  Since \(\delta_i\to0\), this is a contradiction,
proving the unit-scale theorem and, by scaling, Theorem~\ref{thm:chow-iii-perelman-pseudolocality}.
\end{theorem}

\begin{remark}[Role of the scalar-curvature hypothesis]
\label{rem:chow-iii-pseudolocality-scalar-role}
\source{chow-et-al-ricci-flow-part-iii:page-0200}
\dcref{ch21:21.18}
\group{Logarithmic Sobolev contradiction}

The initial lower bound on scalar curvature is used when comparing the localized
entropy with the logarithmic Sobolev inequality.
\end{remark}

\begin{theorem}[Exercise 21.17 (rigidity alternatives for the limiting entropy)]
\label{ex:chow-iii-entropy-rigidity-alternatives}
\dcref{ch21:21.17}
\group{Local entropy}
\source{chow-et-al-ricci-flow-part-iii:page-0196}
\uses{lem:chow-iii-noncollapsed-entropy-limit}
For the limiting nonpositive Harnack quantity, the transition time in the
strong maximum principle is either the terminal time or the initial endpoint
of the rescaled interval.
\end{theorem}
